Con el objetivo de facilitar la comprensión del sistema y organizar las funcionalidades de forma coherente, los requisitos se han agrupado en módulos funcionales que representan las principales áreas de operación del sistema. Estos módulos se describen a continuación de forma general, mientras que en los apartados posteriores se detallan sus requisitos funcionales y de información asociados.

La descomposición del sistema en módulos funcionales responde a los principios de diseño arquitectónico en ingeniería del software, que promueven la organización del sistema en componentes con responsabilidades bien definidas, alta cohesión y bajo acoplamiento, con el objetivo de facilitar su comprensión, mantenimiento y evolución \citep[Ch.~6--7]{sommerville}.

Aunque todos los módulos forman parte del alcance funcional previsto para la solución, su grado de madurez actual no es homogéneo. En el estado actual del proyecto, el módulo M1 constituye el núcleo funcional más consolidado del sistema, mientras que el módulo M2 ya dispone de una base de implementación real asociada a la gestión de clientes, perfiles comerciales, visitas y rutas. El resto de módulos mantienen su validez como estructura de requisitos y como guía para la evolución futura del sistema.

\subsection{M1. Gestión de acceso, alta y autorización de usuarios}

Este módulo agrupa las funcionalidades relacionadas con la identificación de los usuarios del sistema, el alta de nuevas cuentas, la autenticación y el control de acceso según el perfil asignado. Su finalidad es garantizar que únicamente usuarios autorizados puedan utilizar la aplicación y que cada uno acceda exclusivamente a las funciones que le correspondan dentro del sistema.

Asimismo, este módulo contempla la gestión de roles, el tratamiento de solicitudes de registro, la administración del estado de las cuentas, la trazabilidad de accesos y el registro de acciones administrativas asociadas a la gestión de usuarios. Se trata del módulo más transversal de la aplicación, ya que establece los mecanismos de seguridad y control sobre los que se apoya el resto del sistema.

\subsection{M2. Gestión comercial y clientes}

Este módulo agrupa las funcionalidades relacionadas con la gestión de los clientes profesionales y con el seguimiento estructurado de la actividad comercial desarrollada por el distribuidor y sus agentes. Su objetivo es centralizar la información relevante sobre los clientes, facilitar el trabajo comercial diario y permitir un mejor control de las relaciones mantenidas con cada establecimiento.

Durante el desarrollo del proyecto, este módulo ha evolucionado desde una concepción inicial centrada únicamente en la ficha básica del cliente hacia un modelo más completo, en el que también se contempla la existencia de perfiles comerciales, la asignación de clientes a comerciales, la planificación de visitas y la organización de rutas de trabajo. De este modo, el módulo M2 deja de actuar exclusivamente como repositorio de clientes y pasa a constituir la base del soporte operativo a la actividad comercial.

\subsection{M3. Catálogo, productos y coloración}

Este módulo agrupa las funcionalidades relacionadas con la consulta y gestión del catálogo de productos profesionales comercializados por el distribuidor. Su objetivo es proporcionar a los clientes profesionales y a los agentes comerciales una forma estructurada de acceder a la información de los productos, sus características técnicas y su clasificación dentro de las distintas líneas comerciales.

Asimismo, el módulo contempla la organización jerárquica del catálogo mediante categorías y líneas comerciales, así como la gestión de recursos de apoyo, documentación técnica y materiales complementarios asociados a productos o gamas. También incluye herramientas específicas para la consulta de cartas de color y referencias cromáticas utilizadas en los procesos de coloración capilar.

\subsection{M4. Pedidos, entregas y cobros}

Este módulo agrupa las funcionalidades relacionadas con la gestión del proceso comercial de venta de productos, desde la creación de pedidos hasta la entrega de la mercancía y el registro de los cobros asociados. Su objetivo es facilitar tanto a los clientes profesionales como al distribuidor y a los agentes comerciales la gestión de pedidos, el seguimiento de su estado y el control de los pagos realizados o pendientes.

De forma complementaria, el módulo permite gestionar el detalle de cada pedido mediante líneas de productos y conceptos adicionales, así como registrar la información logística asociada a la entrega, incluyendo su modalidad, estado e incidencias. Del mismo modo, contempla el registro de cobros totales o parciales vinculados a pedidos o entregas.

\subsection{M5. Gestión técnica del salón y seguimiento de servicios}

Este módulo agrupa las funcionalidades destinadas a registrar y consultar información técnica sobre los servicios realizados en las peluquerías clientes del sistema. Su objetivo es permitir a los profesionales del salón llevar un seguimiento estructurado de los trabajos realizados a sus clientes, especialmente aquellos relacionados con productos de la marca, facilitando la consulta de historiales técnicos y el análisis del uso de productos.

En este sentido, el módulo contempla la gestión de clientes finales del salón, los servicios o tratamientos realizados, sus fichas técnicas asociadas y los productos empleados en cada intervención. Además, incorpora soporte para la generación de sugerencias de productos y para la elaboración de comunicaciones personalizadas a partir de la información técnica registrada.

\subsection{M6. Promociones, notificaciones y formaciones}

Este módulo agrupa las funcionalidades relacionadas con la comunicación entre el distribuidor y las peluquerías usuarias del sistema. Su objetivo es facilitar la difusión de promociones comerciales, la organización de actividades formativas y el envío de notificaciones o recordatorios relevantes para los usuarios de la aplicación.

Además, este módulo permite segmentar clientes mediante categorías comerciales, asociar promociones a productos, líneas o grupos de clientes y gestionar avisos vinculados a eventos relevantes del sistema. De este modo, no solo cubre la comunicación promocional, sino también la comunicación operativa y la gestión de acciones formativas e inscripciones.

\subsection{M7. Administración e integraciones}

Este módulo agrupa las funcionalidades destinadas a la administración general del sistema y a la integración con herramientas o servicios externos utilizados en la actividad del distribuidor. Su objetivo es permitir la configuración y mantenimiento de la aplicación, facilitar el análisis de la actividad comercial y permitir la interoperabilidad con sistemas externos utilizados en la operativa del negocio.

En particular, este módulo contempla la gestión de parámetros generales de configuración, el registro y seguimiento de integraciones externas y las operaciones de exportación o sincronización de información. Asimismo, incorpora soporte a la generación de propuestas de pedido a proveedor a partir de la información registrada en la plataforma o introducida manualmente.

\subsection{M8. Funcionalidades adicionales}

Este módulo recoge un conjunto de funcionalidades identificadas durante el proceso de análisis de requisitos que, por su naturaleza experimental o complementaria, no forman parte del núcleo principal de la primera versión del sistema. Estas funcionalidades buscan ampliar el valor de la aplicación para los clientes mediante el uso de tecnologías avanzadas y nuevas herramientas de apoyo al trabajo del peluquero.

Entre estas funcionalidades se incluyen, de forma potencial, herramientas de recomendación de tratamientos, simulación visual de coloración, reconocimiento de productos mediante cámara y gestión de reservas online por parte de clientes finales. En consecuencia, estas capacidades podrán incorporarse en fases posteriores del proyecto si el alcance temporal y técnico lo permite.